Anderson, Alan Ross; ed. Minds and Machines. Englewood Cliffs, NJ.: Prentice-Hall, 1964.
 A classic collection of stimulating articles on the mind-body problem, including Alan
 Turing's important paper "Computing Machinery and Intelligence", and J.R. Lucas'
 provocative "Minds, Machines, and Gödel".
Applewhite, Philip. Molecular Gods: How Molecules Determine Our Behavior. Englewood
 Cliffs, NJ.: Prentice-Hall, 1981. A book whose subtitle helped spark my "Careenium"
 dialogue (Chapter 25).
Atlan, Henri. Entre le cristal et la fumee: Essai sur l organisation du vivant. Paris: Editions du
 Seuil, 1979. A biologist who views life as the emergence of ordered complexity out of
 randomness here explains his philosophy.
Axelrod, Robert. The Evolution of Cooperation. New York: Basic Books, 1984. A beautifully
 written account of how mutually beneficial behavior-that is, cooperation-can emerge among
 purely egoistic organisms that share an environment over time.
Ayala, Franciscp Jose and Theodosius Dobzhansky, eds. Studies in the Philosophy of
 Biology: Reduction and Related Problems. Berkeley: University of California Press, 1974.
 The proceedings of one of the most fascinating conferences I have ever heard about. Many
 great biological thinkers were present, and discussed the most fundamental questions about
 how life and minds can be reconciled with physical law. I wish I had been therel
Bandelow, Christoph. Inside Rubik's Cube and Beyond. Boston: Birkhfiuser, 1982- A clear
 and mathematically oriented book about the Cube and its successors. Perhaps the best of all
 the Cube books in English.
Barr, Avron. "Artificial Intelligence: Cognition as Computation". In The Study of
 Information: Interdisciplinary Messages, edited by Fritz hlachlup and Una Mansfield. New
 York: Wiley-Interscience, 1983. A paper putting forward the orthodox Al dogma. that
 mental activity is "information processing"-more specifically, that the manipulation of
 "symbols" (representational data structures) by suitable computer programs is no more and
 no less than what minds do.
Beck, Anatole and David Fowler. "A Pandora's Box of Non-Games". In Seven Years of
 Manifold, . edited by Ian Stewart and John Jaworski. Nantwich, England: Shiva
 Publications, 1981. An amusing collection of silly pseudo-games with more of a moral than
 might first meet the eye.
Beckett, Samuel. Waiting for Godot. New York: Evergreen, 1954. The classic existential
 drama in which meaninglessness pervades-and at whose core is the famous nonsensical
 verbal vomit emitted from the mouth of the sad sack named "Lucky".
Benton, William. Normal Meanings. Paducah, Ky.: Deer Crossing Press, 1978. A book of
 strangely evocative poems, a bit comprehensible in parts, totally incomprehensible in
 others.
Bernstein, Leonard. The Joy of Music. New York: Simon \& Schuster, 1959. An exciting
 medley of ideas by this exuberantly articulate thinker and musician. His dialogues are
 especially enjoyable.
Biggs, John R. Letterforms and Lettering. Poole, England: Blandford Press, 1977. One of the
 best books on the fluidity of letterforms I have run across. Includes short sections on other
 languages, such as Hebrew, Chinese, and Arabic.
Blesser, Barry et al. "Character Recognition based on Phenomenological Attributes". Visible
 Language 7, no. 3 (Summer 1973). An early article by researchers who clearly had come to
 appreciate the depth of the letter-recognition problem.
Bloch, Arthur. Murphy's Law; Murphy's Law, Book Two; Murphy's Law, Book Thres. Los
 Angeles: Price/Stern/Sloan, 1977, 1980, 1982. Humorous and cynical observations about
 the human condition, featuring many self-referential or self-undermining aphorisms.
Boeke, Kees. Cosmic View: The Universe in 40 jumps. New York: John Day, 1957. A book
 to instill humility and awe in anyone, as well as a vivid sense of the meaning of the term
 "astronomical number".


'Bombaugh, Charles Carroll. Oddities and Curiosities of Words And literature, edited
  annotated by Martin Gardner. New York: Dover, 1961. An antique collection of
  palindromes, acrostics, pangrams, and so on, for wordmongers and people who love the
  bizarre fringes of language. (Gardner's footnotes at the back are the best part of the book.)
Boole, George. The Laws of Thought. New York: Dover, 1961. A reprint of the old classic
  from the mid-1850's. The hubris of the title is quite remarkable, especially in light of what
  130 years' progress has revealedl
Brams, Steven, Morton D. Davis, and Philip Strafin, Jr. "The Geometry of the Arms Race".
  International Studies Quarterly 23, no. 4 (December 1979): 567-88. Looking at international
  behavior in terms of the iterated Prisoner's Dilemma and similar payoff matrices.
Brilliant, Ashleigh. I May Not Be Perfect, but Parts of Me Are Excellent: I Have Abandoned
  My Search for Truth, and Am Now Looking for a Good Fantasy; Appreciate Me Now and
  Avoid the Rush ; I Feel Much Better, Now That I've Given Up Hope. Santa Barbara, Calif.:
  Woodbridge Press, 1979-1984. Four books containing many incisive epigrams about life,
  death,, love, relationships, greed, egotism, loneliness, fear, and so on. None is longer than
  seventeen words.
Bush, Donald J. The Streamlined Decade. New York: George Braziller, 1975. Showing how
  the style of an era permeates its creations in all media. Full of elegant photos. Compare with
  the books by Loeb and McCall (see below).
Byrd, Donald. "Music Notation by Computer". Ph.D. thesis, Indiana University Computer -
  Science Department. Bloomington, 1984. About the problems of developing a computer
  program that will have some "understanding" of the subtleties of music notation. The
  program, SMUT, produced the musical examples in this book.
Chaitin, Gregory. "Randomness and Mathematical Proof". Scientific American 232, no. 5
  (May 1975): 47-52. An enlightening way of defining the meaning of "random pattern", and
  the unexpected and deep resonances with GSdel's incompleteness theorem and other
  metamathematicat results.
Charniak, Eugene, C. K. Riesbeck, and Drew V. McDermott. Artifictal Intelligence
  Programming. Hillsdale, NJ.: Lawrence Erlbaum Associates, 1980. Sophisticated ways of
  using Lisp and Lap-like languages in artificial-intelligence research.
Collet, Pierre and Jean-Pierre Eckmann. Iterated Maps on the Interval as Dynamical Systems.
  Boston: Birkhauser, 1980. An in-depth study of the iteration of simple smooth functions on
  the interval [0,1], and the resulting roads to turbulent behavior as parameters are varied.
Compugraphic Corporation. A Portfolio of Text and Display Type. Wilmington, Mass.:
  Compugraphic Corporation, 1982. A collection of many typefaces (mostly book faces),
  including several extensions to other alphabets of faces originally conceived only for our
  alphabet.
Conway, John Horton, Elwyn Berlekamp, and Richard K`Guy. Winning Ways (for your
  mathematical plays).' New York: Academic Press, 1982. A two-volume set of remarkable
  games, some analyzed, some unanalyzed, filled with humorous pictures and an oddball type
  of creative wordplay that is Conway's hallmark. Included in Volume 2 are discussions of the
  Cube and Conway's game of Life.
Coueignoux, Philippe. "La reconnaissance des caractcres". La Recherche 12, no. 126
  (October, 1981): 1094-1103. An interesting article about the workings of some computer
  systems that can read text in a variety of typefaces. The ideas derive from work by Blesser
  and colleagues (see above). This work is rather practically oriented, and does not come
  close to giving computers an understanding of letterforms in their full generality.
Csanyi, Vilmos. General Theory of Evolution. Budapest: Akademiai Kiad6, 1982. A thorough
  investigation of the process of simultaneous evolution on both genetic and "memetic"
  fronts.
Davies, Paul. God and the New Physics. New York: Simon \& Schuster, 1983. In this book,
  Davies tackles the biggies of metaphysics: creation, free will, religion, souls, and more.
  Since he is a good scientist as well as a good writer, his musings are articulate and
  penetrating.


-----Other Worlds New York Simon \& Schuster, 1980 A popular account , by a highly
  reliable professional, of the mysteries at the base of quantum mechanics.
Davis, Morton D. Game Theory: A Nontechnical Introduction. New York: Basic Books,
  1983. An excellent overview of all the main ideas of game theory, including many
  unresolved issues such as the Prisoner's Dilemma.
Dawkins, Richard. The Se\#ish Gene. New York: Oxford University Press, 1976. A book that
  views organisms as by-products of a purely molecular-level competition for efficient self-
  replication. A topsy-turvy, disorienting, yet powerfully revealing viewpoint.
DeLong, Howard. A Profile of Mathematical Logic. Reading, Mass.: Addison-Wesley, 1970.
  A wonderful book on the philosophical and technical issues of logic by someone who
  knows how to achieve an artistic balance between formalisms and ideas that appeal to the
  intuition.
Dennett, Daniel C. Brainstorms: Philosophic Essays on Mind and Psychology. Cambridge,
  Mass.: Bradford Books, MIT Press, 1978. A collection of penetrating analyses of problems
  of mind, brain, and computer models of thought, perception, and sensation. Excellent
  rebuttals of such figures as B. F. Skinner and J. R. Lucas, and a wonderful "dessert" at the
  end. Dennett's graceful style and comparative lack of in-references and jargon make this
  book much more engaging than the average philosophy book.
-----"Can Machines Think?" Unpublished manuscript. Afresh look at the power of the Turing
  Test, completely in sympathy with my view that the test as originally posed is, as valid as
  ever, despite the doubts of many.
-----Cognitive Wheels: The Frame Problem of Artificial Intelligence". In Minds, Machines,
  and Evolution. Edited by C. Hookway. New York: Cambridge University Press, 1985. An
  article about why artificial intelligence is so far from achieving common sense.
-----Elbow Room: The varieties of Free Will Worth Wanting. Cambridge, Mass.: Bradford
  Books, MIT Press, 1984. One provocative metaphor after another, all building up towards
  an image of "self-made selves" that enjoy as much free will as it is reasonable to hanker
  after. Although I resist some of this book's conclusions, I find it the best writing on the
  subject that I know.
-----The Logical Geography of Computational Approaches (A View from the East Pole)".
  Unpublished manuscript. Dennett wittily divides the world of Al approaches into two
  camps: the orthodox one (High Church Computtionalism, centered on the "East Pole"
  (located at MIT), and the unorthodox one ("New Connectionism"), scattered hither and yon.
-----The Myth of the Computer: An Exchange". New York Review of Books. (June 24 1982):
  a civil reply to John Searle´s biting review of The Mind's I.
-----The Self as a Center of Narrative Gravity". In Self and Consciousness, edited by P. M.
  Cole et al. New York: Praeger, 1985. A wonderful new metaphor for thinking about
  abstractions such as "I".
Dewdney, A. K. "Computer Recreations: A computational garden sprouting anagrams,
  pangrams, and few weeds Scientific American 251, no. 4 (October 1984): 20-27. Includes a
  description of Lee Sallows' pangram machine, concluding with a public challenge to
  discover a computer-generated pangram.
DeWitt, Bryce S. and Neill Graham, eds. The Many- Worlds Interpretation of Quantum
  Mechanics Princeton NJ. Princton University Press 1973. A well-rounded presentation of
  one of the most disorienting yet irrefutable ways of thinking about reality yet invented.
Dyer, Michael. In-Depth Understanding. Cambridge, Mass.: MIT Press, 1983. A book
  summarizing a Ph.D. project in language understanding that puts together many recent Al
  ideas about how memory is organized and how flexible control structures fit into the
  picture,,
Dylan, Bob. Tarantula. New York: Macmillan, 1971. The poet-singer lets loose some stream-
  of-consciousness musings that are sometimes reasonably intelligible and sometimes totally
  wacko.
Edson, Russell. The Clam Theater. Middletown, Conn.: Wesleyan University Press, 1973
  Strange and surrealistically written fantasies permeated by a tragic vision of life. To



  me, the most haunting passage is a description of a human head: "this teetering bulb of
  dread and dream ...".
Eidswick, Jack. "How to Solve the n x n X n Cube". Mathematics and Statistics Department,
  University of Nebraska, Lincoln, 1982. A useful manual for those beset by generalized
  cubic frustration.
Endl, Kurt. Rubik's Cube Made Simple; The Pyramid; Pyraminx Cube; Impossiball;
  Megaminx; Rubik's Master Cube. Giessen, Germany: WUrfel-Verlag GmbH, 1982. These
  booklets, together with the previous entry, will give you a good reference shelf on
  Cubology.
Erman, Lee D. et al. "The Hearsay-II Speech-Understanding System: Integrating Knowledge
  to Resolve Uncertainty". ACM Computing Surveys 2, no. 2 (June 1980): 213-53. A review
  article describing, with a good deal of hindsight, the architecture and performance of what I
  consider to be the most inspiring work yet done in artificial intelligence.
Evans, Thomas G. "A Program for the Solution of a Class of Geometric-Analogy Intelligence
  Test Questions". In Semantic Information Processing, edited by Marvin
Minsky. Cambridge, Mass.: MIT Press (1968): 271-353. Perhaps the first truly large system
  ever written in Lisp, this impressive project is often pointed to as having "solved" the
  problem of analogies. How far from the truth that is!
Ewing, John and Czes Koiniowski. Puzzle It Out: Cubes, Groups, and Puzzles. New York:
  Cambridge University Press, 1982. Another lively and mathematically interesting book
  about the Cube and its variants.
Fahlman, Scott E., Geoffrey Hinton, and Terrence J. Sejnowski. "Massively Parallel
  Architectures for Al: NETL, Thistle, and Boltzmann Machines". In "Proceedings of the
  AAAI-83 Conference", Washington, D. C., August, 1983. A comparison of some recent
  ideas for AI architectures, including ones inspired by physics, in which statistical
  emergence plays a central role.
Falletta, Nicholas. The Paradoxicon. New York: Doubleday, 1983. An excellent collection of
  paradoxes of all sorts, ranging from Zeno and Epimenides to modern-day voting paradoxes,
  optical illusions, and antinomies in the philosophies of science and mathematics.
Fauconnier, Gilles. Espaces Mentaux: Aspects de la construction du seas dans les langues
  naturelles. Paris: Les Editions de Minuit, 1984. A probing inquiry into how we make sense
  of frame-crossing and counterfactual statements of this sort: "If Clark Gable had been a
  woman, Scarlett O'Hara would have been a man". Full of ideas about reference, identity,
  slippability, and essence.
Feigenbaum, Mitchell. "Universal Behavior in Nonlinear Systems". Los Alamos Science 1,
  no. 1 (Summer 1981): 4-27. An excellent introduction to the studies of chaos that come
  from iterating simple functions. Excellent graphics, some of which are reproduced in
  Chapter 16.
Feldman, Jerome and Dana Ballard. "Connectionist Models and Their Properties". Cognitive
  Science 6, no. 3 (July-September 1982): 205-54. One of several intriguing approaches to
  computational collectivism currently being investigated in cognitive science. Like most
  such models, this one is strongly influenced by ideas about human perception.
Feynman, Richard P. The Character of Physical Law. Cambridge, Mass.: MIT Press, 1967.
  The transcripts of five marvelous lectures delivered at Cornell University in 1964 by a
  physicist who is not only sharp as a nail but also sparklingly witty.
Feynman, Richard P., Robert B. Leighton, and Matthew Sands. The Feynman Lectures in
  Physics. Reading, Mass.: Addison-Wesley, 1965. Lectures given to beginning students,
  often more suitable for advanced students, filled with the excitement of science and the
  peppery observations of Feynman.
Franke, H. W. Computer Graphics-Computer Art. New York: Phaidon, 1971. By far the best
  collection of examples and discussion of computers and art I have yet seen, even though it
  is quite old by now.
Frey, Alexander H., Jr. and David Singmaster. Handbook of Cubik Math. Hillside, NJ.:
  Enslow, 1982. The group-theoretical ideas behind the cube are developed systematically,
  for possible use in an advanced high-school or college-level course.


Friedman, Daniel P. The Little LISPer. Chicago: Science Research Associates, 1974. an
  engaging and elegant introduction to the recursive ideas of Lisp, by one of its most ardent
  and accomplished practitioners.
Fromkin, Victoria A, ed. Errors in Linguistic Performance: Slips of the Tongue, Ear, Pen, and
  Hand. New York: Academic Press, 1980. A compendium of articles on how we can infer
  general mechanisms of thought from observing everyday surface-level quirks that show up
  ubiquitously in speech, typing, handwriting, listening, and so on.
Frutiger, Adrian. Type Sign Symbol. Zurich: Editions ABC, 1980. This book by the creator of
  Frutiger and Univers, two of today's most popular and elegant sans-serif typefaces, deals
  with the interaction of technical constraints and artistic creation, and features many
  beautiful letterforms and styles.
Gablik, Suzi. Progress in Art. New York: Rizzoli International, 1976. A heretical theory
  suggesting that there is a reason why art is where it is today, and that art follows a
  comprehensible trajectory in some abstract space, even if it is explicable only after the fact.
Gardner, Martin. Fads and Fallacies. New York: Dover, 1952. A book that comes about as
  close as one could hope to being a course on common sense. Many chapters of this classic
  bunk-puncturer should be part of the cultural education of everyone.
Gardner, Martin. "Mathematical Games: White and brown music, fractal curves, and one-
  over-f fluctuations". Scientific American 238, no. 4 (April 1978): 16-32. A solid
  introduction to the concepts of fractals and multi-level statistical structures.
Gardner, Martin. Science: Good, Bad, and Bogus. Buffalo: Prometheus Press, 1981.
  Following in the footsteps of Fads and Fallacies, this book sharp-swordedly decapitates one
  dragon of hokum after another-and yet sadly, like a hydra's many heads, they always seem
  to pop back up again.
Gardner, Martin. Wheels, Life, and Other Mathematical Amusements. New York: W. H.
  Freeman, 1983. Gardner's final collection of his sparkling Scientific American columns,
  showing why mathematics is the ultimate metamagic.
Gebstadter, Egbert B. Thetamagical Memas: Seeking the Whence of Letter and Spirit. Perth:
  Acidic Books, 1985. A curious pot-pourri, bloated and muddled-yet remarkably similar to
  the present work. This is a collection of Gebstadter's monthly rows in Literary Australian
  together with a few other articles, all with prescripts. Gebstadter is well known for his love
  of twisty analogies, such as this one (unfortunately not found in his book): "Egbert
  Gebstadter is the Egbert Gebstadter of indirect self-reference."
Geisel, T. and J. Nierwetberg. "Universal Fine Structure of the Chaotic Region in Period-
  Doubling Systems". Physical Review Letters 47, no. 14 (October 5, 1981): 975-78. An
  investigation turning up remarkable order deep in the heart of chaos.
Gentner, Dedre. "Structure-Mapping: A Theoretical Framework for Analogy". Cognitive
  Science 7 (1983): 155-70. One in a series of papers attempting to draw guidelines that will
  help distinguish good analogies from bad.
Gödel, Kurt. On Formally Undecidable Propositions in "Principia Mathematica" and Related
  Systems. Translated by Bernard Meltzer, edited and with an introduction by R. B.
  Braithwaite. New York: Basic Books, 1962. The fundamental work that opened up worlds
  to logicians, philosophers, computer scientists, and others.
Golden, Michael. "Don't Rewrite the Bible". Newsweek (November 7,4983):47. This piece
  sounds almost like my "Person Paper" in the way it exploits crude mockery of progressive
  new usages in order to make oppressive old usages sound good. The difference is that
  Golden is not being sarcastic.
Golomb, Solomon. "Rubik's Cube and Quarks". American Scientist 70, no. 3 (May June
  1982): 257-59. In which Golomb puts forth his analogy between twisted corners on the
  Cube and fractionally charged subatomic particles.
Gonick, Larry and Mark Wheelis. A Cartoon Guide to Genetics. New York: Barnes \& Noble,
  1983. A short and snappy-but most informative and entertaining-overview of molecular
  biology, genetics, and their human import.
Gould, Stephen Jay. Hen's Teeth and Horse's Toes. New York: W. W. Norton, 1983. Written
  in a lively and engaging way, this book (like its predecessors, Ever Since Darwin and


 The Panda's Thumb) relates with great clarity the issues of evolution as illustrated in
 nature's unbelievable variety of quirks.
Gray, J. Patrick and Linda Wolfe. "The Loving Parent Meets the Selfish Gene". Inquiry 23:
 233-42. A cogent set of arguments against the pessimistic theses of Louis Pascal, to which
 he then replies (see below).
Grossman, I. and W. Magnus. Groups and Their Graphs. New York: Random House New
 Mathematical Library, 1975. An excellent visual introduction to group theory, featuring so-
 called "Cayley diagrams", which lay bare the structure of a group in a marvelously clear
 way.
Ground Zero. Nuclear War: What's in It for You? New York: Pocket Books, 1982. A
 dispassionate, calmly reasoned discussion for "just-plain folk" of what it means for
 countries to make and threaten each other with nuclear weapons. The best primer that I have
 come across on the subject, and one that I wish everyone would read and take to heart.
Guillemin, Victor. The Story of Quantum Mechanics. New York: Scribner's, 1968. A
 thoughtful survey of where quantum mechanics came from and what it has revealed about
 the physical world, concluding with a long discussion of the philosophical implications of
 quantum mechanics for such things as free will and causality.
Haab, A,, A. Stocker, and W. Hattenschweiler. Lettera 1 and Lettera 2. Arthur Niggli, 1954
 and 1961. Some deliciously fanciful alphabets are featured in these books. Books like these
 ought to make anyone marvel at the fluidity of letters and human perception.
Hachtman, Tom. Double Takes. New York: Harmony Books, 1984. Two-in-one caricatures,
 done in a very satisfying way. It makes one wonder: could a computer ever do this kind of
 thing?
Hansel, C. E. M. ESP and Parapsychology: A Critical Re-Evaluation. Buffalo: Prometheus
 Books, 1980. A scholar looks at the claims of parapsychologists and finds them wanting.
 Hanson, Norwood Russell. Patterns of Discovery. New York: Cambridge University Press,
 1969. A philosophically rich discussion of how physical concepts progressed by fits and
 starts through the centuries, culminating in the strange world of elementary particles that we
 are now finding are not so elementary.
Hardin, Garrett. "The Tragedy of the Commons". Science 162, no. 3859 (December 13,
 1968): 1243-48. Using the metaphor of shared grazing land that gets ruined by overuse
 without any specific person being responsible, Hardin argues with great force that people's
 narrow-minded selfishness will drive us to extinction unless we band together and form
 strong organizations dedicated to global ecological goals, particularly population control.
Harel, David. "Response to Scherlis and Wolper". Communications of the ACM 23, no. 12
 (December 1980): 736-37. An amusing list of tricky self-references known to Harel,
 sporting a footnote claiming to be the first published example of "a reference to the very
 self-reference being discussed in the very letter being read!"
Hart, H. L. A. "Self-Referring Laws". In Festskrift Tillagnad Karl Olivecrona. Stockholm:
 Kungliga Boktryckeriet, P. A. Norstedt och Soner, 1964. A reply to Alf Ross (see below), in
 which Hart, perhaps today's leading philosopher of law, claims that self-amending laws are
 legally possible.
Harth, Erich. Windows on the Mind: Reflections' on the Physical Basis of Consciousness.
 New York: William Morrow, 1982. Lively discussions of brain, mind, consciousness, and
 how they all tie together, by a physicist with a philosophical bent.
Haugeland, John, ed. Mind Design: Philosophy, Psychology, and Artificial Intelligence.
 Cambridge, Mass.: Bradford Books, MIT Press, 1981. A self-proclaimed sequel to Alan
 Ross Anderson's Minds and Machines, now that artificial intelligence has had a couple of
 decades to develop. This volume is of uneven quality, but it contains many articles sure to
 provoke thought.
Heiser, Jon F. et al. "Can Psychiatrists Distinguish a Computer Program Simulation of
 Paranoia from the Real Thing? The Limitations of Turing-Like Tests as Measures of the
 Adequacy of Simulations". Journal of Psychiatric Research 15, no. 3 (1979): 149-62.
 Researchers who developed the program called "Parry" claim that the Turing Test's validity
 is cast in doubt because Parry passed a pseudo-Turing Test.


Hinton, Geoffrey and James Anderson, eds. Parallel Models of Associative Memory.
  Hillsdale, N .J.: Lawrence Erlbaum Associates, 1981. A collection of articles exploring
  different models, all of which are based on the thesis that human perception and memory
  are statistically emergent phenomena and are best modeled as such.
Hintze, Wolfgang. Der Ungarische Zauberwiirfel. Berlin: VEB Deutscher Verlag der
  Wissenschaften, 1982. This is an excellent book about the Cube, including new
  mathematical results on subgroups. Serious cubists who can read German should obtain this
  book.
Hobby, John, and Gu Guoan. "A Chinese Meta-Font". Stanford, Calif.: Stanford University
  Computer Science Department Technical Report STAN-CS-83-974, 1983. A description of
  a system for creating Chinese characters whose style is "tunable" to some extent. This
  system is very similar in some ways to Han Zi, the variable-style Chinese-character-
  generation system developed by David Leake and myself at Indiana, at about the same time.
Hodges, Andrew. Alan Turing: The Enigma. New York: Simon \& Schuster, 1983. The
  definitive biography of this extremely important figure in mathematics and computer
  science, written with empathy, insight, and charm.
Hofstadter, Douglas R. Ambigrams. Forthcoming in 1985 from the Centre d'Art
  Contemporain (Geneva, Switzerland). A collection of ambigrams (or "inversions", as Scott
  Kim prefers to call them) aptly subtitled "A Panoply of Palindromic Pinwheels". The
  foreword is by Scott (to whose book I symmetrically wrote the foreword). Also included are
  a lecture on ambigrams (in French) and an interview (in English).
-----Analogies and Metaphors to Explain Godel's Theorem". Two-Year College Mathematics
  Journal 13, no. 2 (March, 1982): 98-114. A collection of simple images and ideas that help
  to build up one's intuitions to the point where one can fully digest the intricacies of Godel's
  clever construction.
-----The Architecture of Jumbo". In "Proceedings of the Second Machine Learning
  Workshop", Monticello, Illinois, 1983. A description of the architecture of a biologically-
  inspired Al system that exploits parallelism and randomness in order to build up "well-
  chunked wholes" out of isolated parts and then to allow them to seek maximal "happiness"
  by internally reconfiguring themselves in a manner governed by a computational
  temperature.
-----The Copycat Project: An Experiment in Nondeterminism and Creative Analogies".
  Cambridge, Mass.: MIT Artificial Intelligence Laboratory Al Memo 755, April 1984. A
  description of ongoing research in my group, first at Indiana, then at MIT, and now at
  Michigan. The goal is to make a system capable of doing analogies in a tiny domain, but
  with the insight (and short sight) of humans. Describing how the jumbo architecture can be
  adapted to this task is the burden of this paper. .
-----512 Words on Recursion". Math Bulletin, Bronx High School of Science (1984): 18-19.
  A recursively structured article, which quotes a compressed version of itself (in which, of
  course, a more compressed version is quoted, etc.).
-----Gödel, Escher, Bach: an Eternal Golden Braid. New York: Basic Books, 1979. Originally
  titled "Gödel’s Theorem and the Human Brain", this book tries to build up the concepts
  necessary to show how ideas that arise in Gödel’s proof will one day be at the core of
  explanations of consciousness and the meaning of the word "I". Analogy, humor, and
  contrapuntal dialogues are essential features of GEB.
-----Gridfonts". Unpublished manuscript, 1984. A still-growing collection of (currently) about
  300 skeletal typefaces-that is, distinctive and artistically consistent ways of rendering all 26
  letters in a highly confined grid. The purpose is to provide material for discussion of the
  roots of style and the fluid nature of categories.
-----On Seeking Whence". Unpublished manuscript, 1982. A discussion of the immense
  difficulty of extrapolating linear patterns, and its relationship to such deeply human
  experiences as scientific induction, analogy-making, understanding music, and creativity.
-----Poland: A Quest for Personal Meaning". Poland (May, 1981): 42-47. A compression of a
  longer piece called "Poland: A Mythical Quest", describing my powerful emotional
  reactions to Poland on my first visit there, in 1975.


-----In Search of Essence". Unpublished manuscript, 1983. A discussion of what the "essence"
  of a written passage is, by way of a description of problems encountered in translating the
  dialogue Controcrostipunctus (from C,6, Escher, Bach) into French.
-----Simple and Not-So-Simple Analogies in the Copycat Domain". Unpublished manuscript.
  A collection of 84 analogy problems in the Copycat domain that show how diverse the tiny
  alphabetic universe can be.
Hofstadter, Douglas R., Gray Clossman, and Marsha Meredith. "Shakespeare's Plays Weren't
  Written by Him, but by Someone Else of the Same Name". Bloomington: Indiana
  University Computer Science Department Technical Report 96, 1980. "A study of
  intensionality in frame-based representation systems" is the subtitle. The paper's main
  purpose is to discuss the relationship between slots that an entity fills, and that entity's
  identity. The notion of "Core ID" is presented and explored in an analysis of the title
  sentence. Closely related to the problems discussed by Fauconnier (see above).
Hofstadter, Douglas R. and Daniel C. Dennett, eds. The Mind's : Fantasies and Reflections on
  Self and Soul. New York: Basic Books, 1981. Our purpose was to jolt people on all sides of
  the fence in this anthology of fictional and evocative pieces on the curious fact (or illusion)
  that something we call an "I" is somehow connected to some hunk of matter floating
  somewhere and somewhen in some universe.
Holland, John et al. Induction: Processes of Inference, Learning, and Discovery. Forthcoming.
  A wide-ranging study of how learning takes place in genuine human minds, and how
  aspects of it might be modeled in computer simulations. This pioneering work is by a
  computer scientist, two' psychologists, and a philosopher, and its breadth reflects theirr
  diversity.
Hollis, Martin and Steven Lukes, eds. Rationality and Relativism. Cambridge, Mass.: MIT
  Press, 1982. Showing how scholars can get mired in the endlessly circular reasoning that
  arises when any system of belief tries to provide its own justification. Reminiscent of
  "MUnchhausens Zopf": the quandary of Baron MUnchhausen trying to pull himself out of a
  quicksand quagmire simply by tugging on his own braid.
Hopfield, J. J. "Neural networks and physical systems with emergent computational abilities".
  In Proceedings of the National Academy of Sciences 79, Washington, D.C., 1982: 2554-58.
  A system that makes use of statistics and parallelism to model familiar properties of the
  mind: reconstructing a full memory from a partial one, generalizing, and so on.
How to Learn Lettering. Hong Kong: Nam San Publisher, n.d. A collection of Chinese
  characters in modern artistic styles, mostly for the use of advertisers, but also a treasure
  trove for those interested in the elusive "sameness" we see in wildly different renderings of
  a single Platonic essence.
Huff, William. "A Catalogue of Parquet Deformations". School of Architecture, State
  University of New York at Buffalo. Approximately 35 parquet deformations created by
  students in William Hull's studio, assembled and commented on by Huff.
Hughes, Patrick and George Brecht. Vicious Circles and Infinity: An Anthology of
  Paradoxes. New York: Penguin, 1975. A' delightful and provocative collection of
  paradoxical material and choice epigrams embodying paradoxes of every conceivable sort.
  Includes lengthy discussions of the paradox of the Unexpected Examination, also known as
  the "Hangman's Paradox".
Huneker, James Gibbons. Chopin: The Man and His Music. New York: Dover, 1966.
  Originally published around the turn of the century, this romantic biography features florid
  prose so rich that it verges on meaninglessness, and yet it is wonderfully evocative.
Jaspert, W. Pincus, W. Turner Berry, and A. F. Johnson. The Encyclopaedia of Type Faces.
  Poole, England: Blandford Press, 1983. An engrossing catalogue of over 1,000 typefaces.
  One of its unique features is its thorough crediting of designers: the people whose
  exquisitely developed sense of line and space is far too often completely taken for granted.
  An index to designers and an index to typefaces are included. This way you can get a sense
  for the stylistic range of various designers.
Johnson-L'Laird, P. N. and P. C. Wason, eds. Thinking. New' York: Cambridge University
  Press, 1975. A diverse collection of pieces by top-notch authors about nearly


  all aspects of human thought, including imagery, perception, inference, categories, memory,
  language, and so on.
Kadanoff, Leo. "Roads to Chaos". Physics Today 36, no. 12 (December 1983): 46-53. A good
  summary of the connection between mathematical approaches to chaos and the physical
  phenomena to be explained.
Kamack, H. J. and T. R. Keane. "The Rubik Tesseract". Unpublished manuscript, 1982. A
  companion piece to the one by Eidswick (see above), this discussion of a 3 x 3 x 3 x 3
  "Rubik's hypercube" is a tour de force in visualization-for its readers almost as much as for
  its authors!
Kahneman, Daniel and Amos Tversky. "The Simulation Heuristic". In Judgment Under
  Uncertainty: Heuristics and Biases, edited by Daniel Kahneman, P. Slovic, and A. Tversky.
  New York: Cambridge University Press, 1982: 202-8. In what ways are people inclined to
  let situations mentally glide into alternate versions of themselves, and how do subtle
  cognitive pressures modify those tendencies? Two outstanding cognitive psychologists here
  describe their findings about this sort of slippability or "alternity" in their subjects' minds.
Kanerva, Pentti. Self-Propagating Search: A Unified Theory of Memory. Stanford, Calif.:
  Center for the Study of Language and Information, Technical Report, Stanford
University, 1984. A neurocomputational theory of memory: How to make similarity-sensitive,
  reconstructive software out of statistically robust addressing hardware. This may sound a bit
  abstruse, but in my opinion Kanerva's elegant work constitutes one of the most important
  steps toward reconciliation of fluidity and mechanism-mind and brain -that I have seen
  taken.
Kennedy, Paul E. Modern Display Alphabets. New York: Dover, 1974. An excellent
  collection of visually attractive letterforms in styles of all sorts, and exhibiting all degrees of
  wildness.
Kim, Scott E. "The Impossible Skew Quadrilateral: A' Four-Dimensional Optical Illusion". In
  Proceedings of the 1978 AAAS Symposium on Hypergraphics: Visualizing Complex
  Relationships in Art and Science. Boulder, Colo.: Westview Press, 1978. Like the work by
  Kamack and Keane (see above), this article takes a powerful visual imagination to write or
  to read. The familiar "impossible triangle" is here carried one step further, by a bold process
  of analogy. The writing style is full of typical Kimian recursive trickery and parallel
  passages.
-----Inversions. Peterborough, N.H.: Byte Books, 1981. A collection of inversions (or
  "ambigrams", as I prefer to call them) aptly subtitled "A Catalogue of Calligraphic
  Cartwheels". The foreword is by myself (to whose book Scott symmetrically wrote the
  foreword). The prose is as full of illusions and tricks of parallelism as the drawings are,
  although it is not immediately apparent.
-----Noneuclidean Harmony". In The Mathematical Gardner, edited by David A. Klarner.
  Belmont, Calif.: Wadsworth International/Prindle, Weber, and Schmidt, 1981. A serious
  treatise on atonal geometry masquerading as a piece of humor. Describes Georg Cantor's
  important results on supersonic pitches, 2 against 3 correspondence, and uncountable
  rhythms.
-----Visual Art: The Creative Cycle". Response, no. I (December, 1981). Minnesota Artists
  Exhibition Program. A short article putting forth a thesis about creative activity, at the heart
  of which is a loop similar to the one I describe in the P.S. to Chapter 12 of this book.
Kirkpatrick, S., C. D. Gelatt, Jr., and M. P. Vecchi. "Optimization by Simulated Annealing".
  Science 220, no. 4598 (May 13, 1983): 671-80. An article describing a new statistical
  technique for seeking the optimal state of a system with a huge number of degrees of
  freedom, based on local improvement modulated by occasional local degradation, the
  amount of which is determined by a "cooling schedule", analogous to the annealing of a
  metal to strengthen it.
Kleppner, Daniel, Michael G. Littman, and Myron L. Zimmerman. "Highly Excited Atoms".
  Scientific American 244, no. 5 (May 1981): 130-49. A bridge between microscopic
  quantum-mechanical phenomena and macroscopic classical phenomena is provided by the
  study of "Rydberg atoms", with which this paper is concerned.


Knuth, Donald. "The Concept of a Meta-Font". Visible Language 16, no. 1 (Winter 1982): 3-
  27. Knuth describes his program, METAFONT, for creating an entire family of typefaces at
  one fell swoop, by parametrizing all letters of the alphabet with a consistent set of
  parameters. This is the paper that triggered my response printed as Chapter 13.
-----TEX and Metafont: New Directions in Typesetting. Bedford, Mass.: Digital Press, 1979.
  Knuth describes in detail how to use his programs that facilitate typesetting and letter
  design.
Kobylariska, Krystyna. Chopin in His Own Land. Cracow: Polish Music Publishers, 1956. A
  collection of hard-to-find Chopiniana, beautifully reproduced in large format.
Koestler, Arthur. The Act of Creation. New York: Dell, 1964. The great novelist and amateur
  psychologist delves into his own mind to come up with interesting but often wrong-headed
  speculations on humor, creativity, and insight.
-----The Roots of Coincidence: An Excursion into Parapsychology. New York: Vintage,
  1972. Koestler reveals his belief in the occult, somewhat reminiscent of Arthur Conan
  Doyle's belief in fairies.
Kolata, Gina. "Does Godel's Theorem Matter to Mathematics?" Science 218 (November 19,
  1982): 779-80. An excellent description of how the specification of what today would be
  considered "large integers" depends on abstruse concepts from mathematical logic.
Koning, H. and J. Eizenberg. "The language of the prairie: Frank Lloyd Wright's prairie
  houses". Environment and Planning B, 8 (1981): 295-323. A description of how pseudo-
  Frank-Lloyd-Wright houses can be generated from a "shape grammar". Kripke, Saul.
  Naming and Necessity. Cambridge, Mass.: Harvard University Press, 1972. This theory of
  extensions, intensions, "rigid designators", and identity is. in direct opposition to my views
  (put forward in the "Shakespeare" paper cited above) that the roots of identity come from
  patterns of embeddedness in a network.
Kuwayama, Yasaburo. Trademarks and Symbols, Volume 1: Alphabetical Designs. New
  York: Van Nostrand Reinhold, 1973. A celebration of the craziness of letterforms, this
  collection contains letters to boggle anyone's mind.
Larcher, Jean. Fantastic Alphabets. New York: Dover, 1976. The letterforms in Kuwayama's
  book are highly imaginative, but by comparison, the letterforms in Larcher's book are
  downright zany. It just goes to show that a computer program that could deal with letters in
  their full complexity would be able to deal with the entire world.
Lehninger, Albert. Biochemistry, 2d. ed. New York: Worth, 1975. A well-written treatise
  covering the huge field of biochemistry in a lucid manner.
Lennon, John. In His Own Write and A Spaniard in the Works. New York: Simon \&
  Schuster, 1964 and 1965. This Beatle had a wonderful sense of silliness, all his own.
  Stories, pictures, poems-even a little drama or two. Perhaps better than the Beatles' music, if
  I may venture a heretical opinion.
Letraset, Inc. Graphic Art Materials Reference Manual. Paramus, NJ.: Letraset, Inc., 1981.
  This is the best inexpensive typeface manual available, and is well worth the price. It
  contains nearly all the common typefaces, as well as many unusual ones. This is a great way
  to get into typefaces.
Letraset, Inc. Letraset Greek Series. Athens: A. Pallis, 1984. This to me is an amazing
  catalogue of typefaces, because each and every one of them represents a "transalphabetic
  leap". Here we see such faces as Blanchard, Futura Black, Helvetica, Korinna, Optima,
  Souvenir, University Roman, Zipper, and a number of others-all in Greek! All I can say is,
  "Wow!"
Levin, Michael. "Mathematical Logic for Computer Scientists". Cambridge, Mass.: MIT
  Laboratory for Computer Science Technical Report LCS TR 131, June, 1974. An
  unorthodox treatment of logic for those more comfortable with Lisp than with number
  theory.
Lipman, Jean and Richard Marshall. Art about Art. New York: E. P. Dutton, 1978. An
  amusing annotated collection of self-conscious art, including pieces by Roy Lichtenstein,
  Andy Warhol, Robert Rauschenberg, Jasper Johns, Robert Arneson, Tom Wesselmann,
  Larry Rivers, Mel Ramos, Peter Saul, and many others.


Loeb, Marcia. New Art Deco Alphabets. New York: Dover, 1975. Showing how somebody
  can-perfectly recreate the spirit of an era in many different ways. These original alphabets
  all look like they came straight out of the 1930's. Compare with the books by Bush (see
  above), and McCall (see below).
Lucas, J. R. "Minds, Machines, and Gbdel". Reprinted in Minds and Machines, edited by
  Alan Ross Anderson. Englewood Cliffs, NJ.: Prentice-Hall, 1964. The article that launched
  a thousand rebuttals. Although I find its conclusions totally unjustified, the issues are well
  raised and deserve to be thought about far more. Closely related to the issues discussed in
  the book by Webb (see below).
Machlup, Fritz and Una Mansfield, eds. The Study of Information: Interdisciplinary
  Messages. New York: Wiley-Interscience, 1983. A multi-faceted collection of articles by
  high-level scholars about information theory, cybernetics, artificial intelligence, cognitive
  science, librarianship, and more. Many of the pieces are statements of opinion, and conflict
  with one another. This volume includes the exchange between Allen Newell and myself,
  which is continued in the P.S. to my Chapter 26.
Mandelbrot, Benoit. The Fractal Geometry of Nature. New York: W. H. Freeman, 1982. This
  latest presentation of Mandelbrot's visions is rich in imagery and ideas. Fractal shapes
  constitute one of the twentieth century's most fertile mathematical playgrounds.
Marek, George R. and Maria Gordon-Smith. Chopin. New York: Harper \& Row, 1978. One
  of the many biographies of Chopin in English. I find it to be balanced and well written.
Marx, George, Eva GajzSgo, and Peter Gnadig. "The universe of Rubik's cube". European
  journal of Physics 3 (1982): 34-43. The implications of extending the concept of entropy to
  the surface of the Cube are reported on, with the results of several statistical studies done on
  computer.
Max, Nelson. Space-Filling Curves. Chicago: International Film Bureau, Inc., 1974. A
  marvelous movie done with computer graphics showing the fascination of some of the
  earliest-discovered fractal curves. The computer-produced music adds extra appeal and
  eerieness.
May, Robert. "Simple Mathematical Models with Very Complicated Dynamics". Nature 261,
  no. 5560 (June 10, 1976): 459-67. An early review article about the chaos one is led to
  when one iterates simple functions on the interval [0,1 ]. Very informative, except for the
  reversal of two figures!
Maynard-Smith, John. Evolution and the Theory of Games. New York: Cambridge University
  Press, 1982. Showing how new depths of understanding of the mechanisms of evolution can
  come from modeling the interactions of organisms in a common environment in terms of
  mathematical game theory.
McCall, Bruce. Zany Afternoons. New York: Alfred A. Knopf, 1982. A genius for capturing
  the flavor of things turns all his talents to recreating the twenties, thirties, forties, and fifties-
  not only in America but also abroad. This very funny book is an amazing achievement.
  Compare with the books by Bush and Loeb (see above).
McCarthy, John. "Attributing Mental Qualities to Machines". In Philosophical Perspectives
  on Artificial Intelligence, edited by Martin Ringle, 161-95. Atlantic Highlands, N. J.:
  Humanities Press, 1979. Under what circumstances does it make sense to speak of machines
  having desires or beliefs? When to adopt this "intentional stance" is here discussed by one
  of the pioneers of artificial intelligence.
-----History of Lisp". In History of Programming Languages, edited by Richard L.
  Wexelblatt, 217-23. New York: Academic Press, 1980. Lisp's genesis as recalled by its
  inventor.
McCarty, L. T. and N. S. Sridharan. "A Computational Theory of Legal Argument". New
  Brunswick, NJ.: Rutgers University Computer Science Department Technical Report LPR-
  TR-13, January 1982. A professor of law and a computer scientist work together to
  implement "prototype deformation", which they consider the key to computer modeling of
  argumentation by precedent.
McClelland, James L., David E. Rumelhart, and Geoffrey E. Hinton, eds. Interactive
  Activation: A Framework for Information Processing. Forthcoming. A collection of articles


  on statistically emergent parallel computation featuring "temperature" as a regulator of
  randomness.
McKay, Michael D. and Michael S. Waterman. "Self-Descriptive Strings". Mathematical
  Gazette 66, no. 435 (1982): 1-4. A mathematical study of a simple class of sentences that
  attempt to inventory themselves.
Meehan, James R. "Tale-Spin, an Interactive Program that Writes Stories". In Proceedings of
  the 5th International Conference on Artificial Intelligence-1977, Vol. 1: 91-98. Pittsburgh:
  Department of Computer Science, Carnegie-Mellon University. One of the funniest (and
  therefore most sensible) articles ever written on computers and their "understanding" of
  language. Meehan explains how, in one "mis-spun tale", his program made gravity fall into
  a river and drown. Meehan's discussions of mis-spun tales in general are among the most
  illuminating passages I know on that perennial vexation: the fact that AI programs
  stubbornly resist acquiring common sense, no matter how hard or how often they are
  spanked.
Miller, Casey and Kate Swift. The Handbook of Nonsexist Writing for Writers, Editors, and
  Speakers. New York: Barnes and Noble, 1980. A first-class book. It is a scandal that this
  book is not found on every newsman's desk. Thinking about these issues is not only socially
  important, but challenging and fascinating.
Words and Women. New York: Doubleday, Anchor Press, 1977. A magnificent and
  devastating job of showing the absurdity of thinking that women have been dealt a fair hand
  by our society. If language is the thermometer of society, then this book reveals that
  collectively, we are gravely ill.
Mills, George. "Gddel's Theorem and the Existence of Large Numbers". Northfield,
  Minnesota: Mathematics Department, Carleton College, October 1981. Showing how the
  descriptions of some stupendously large numbers can be obtained via recursive definitions
  of functions and the process of diagonalization. Repeated jootsing leads to some remarkable
  destinations!
Minsky, Marvin. "A Framework for Representing Knowledge". In The Psychology of
  Computer Vision, edited by P. H. Winston. New York: McGraw-Hill, 1975. The paper that
  defined such now-central notions to Al as frames, slots, and default assumptions.
-----Matter, Mind, and Models". In Semantic Information Processing, edited by Marvin
  Minsky. Cambridge, Mass.: MIT Press, 1968. Some cogent speculations about the meaning
  of the word "I" in computational terms.
-----Why People Think Computers Can't". AI Magazine (Fall 1982): 3-15. With his usual
  unusual insight, Minsky tears people apart for not understanding how to think about
  thinking (or about machines, for that matter).
Mondrian, Piet. Tout l'tEuvre Peint de Piet Mondrian. Paris: Flammarion, 1976. For a grand
  overview of the evolution of the style of one painter, I know of no better book than this,
  which traces Mondrian from his earliest representational paintings to his most abstract and
  geometrical ones, revealing the sweep to be continuous and logical, but no less dramatic for
  that.
Morrison, Philip and Phylis, and the Office of Charles and Ray Eames. Powers of Ten. New
  York: Scientific American Books, 1983. Kees Boeke's inspirational Cosmic View (see
  above) revisited some thirty years later, with considerably more commentary. A charming
  and . - worthy successor.
Myhill, John. "Some Philosophical Implications of Mathematical Logic: Three Classes of
  Ideas". Review of Metaphysics 6, no. 2 (December 1952): 165-98. In 1982 I met Myhill and
  asked him about this paper. He told me he considered it to be a piece of junk. That
  astonished me, since I consider it to be a thoughtful and important piece of philosophizing
  making use of mathematical metaphors, something that hardly anyone dares to do. You
  never know how someone will evaluate their own work 30 years later!
Nagel, Ernest, and J. R. Newman. Gödel’s Proof. New York: New York University Press,
  1958. A gracious and highly accessible introduction to the twists in Godel's reasoning, as
  well as to the philosophical issues surrounding his work.



Nakanishi, Akira. Writing Systems of the World. Rutland, Vt.: Tuttle, 1980. For a sampling
  of the many different spirits residing in letterforms, see this book. It features reproductions
  of newspaper pages showing each writing system in several different styles.
Newell, Allen. "Physical Symbol Systems". Cognitive Science 4, no. 2 (April-June 1980):
  135-83. A lengthy article putting forth the orthodox dogma on which artificial intelligence
  has traditionally considered itself founded: the idea that a universal computer embodies all
  the prerequisites for intelligent behavior.
Norman, Donald. "Categorization of Action Slips". Psychology Review 88, no. 1 (January
  1981): 1-15. A delightful compendium of types of error that people commit in performing
  everyday actions such as putting water on to boil, answering the telephone, unbuckling
  one's seatbelt, driving home from work, and so on. There is remarkable regularity behind
  the seeming chaos, and Norman's purpose is to chart and exploit that regularity in order to
  reveal hidden mechanisms of thought.
Nozick, Robert. Philosophical Explanations. Cambridge, Mass.: Harvard University, Belknap
  Press, 1981. A very wide-ranging book on philosophy, intended for lay readers as well as
  for professionals. It covers matters from personal identity and the meaning of reference to
  free will and morality, and only occasionally lapses into brief spasms of absolutely
  inscrutable jargon.
Parfit, Derek. Reasons and Persons. Oxford: Clarendon Press, 1984. A very thoughtful
  treatise on moral dilemmas and ethical behavior, rooted in close consideration of the
  deepest roots of caring: why do I care about my present and future self? Parfit knows that to
  make any serious attempt to answer this riddle, one must look long and hard at the meaning
  of the word "I" in the real world and in many counterfactual ones, which he does with skill
  and insight.
Pascal, Louis. "Human Tragedy and Natural Selection". Inquiry 21: 443-60. Taking up where
  Garrett Hardin left off (see above), Pascal paints a gloomy picture of a population explosion
  as the natural outcome of selection itself, something built into the nature of societies just as
  deeply as the sexual drive is built into individuals.
"Rejoinder to Gray and Wolfe". Inquiry 23: 242-51. A bitter indictment of the human race's
  apathy before visibly onrushing catastrophe.
Peattie, Lisa. "Normalizing the Unthinkable". Bulletin of the Atomic Scientists 40, no. 3
  (March 1984): 32-36. Like Pascal, Peattie is concerned with people's apparent ability to turn
  off their sensitivities and to focus on the very local to such an extent that great tragedies are
  allowed to ensue. In a striking analogy, she likens the current public apathy about the arms
  madness to the inhumanity of collaborators in Hitler's concentration camps.
Perec, Georges. La Disparition. Paris: Editions Denoel, 1969. If anything, writing without 'e's
  is harder in French than in English, yet here is an entire novel in that bizarre dialect.
  Naturally enough, its subject is the mysterious disappearance of item number five in a
  collection of twenty-six objects. It was probably inspired by the 'e'-less novel Gadsby,
  written in English by Ernest Vincent Wright in the late 1930's.
Perfect, Christopher and Gordon Rookledge. Rookledge's International Typefinder. New
  York: Frederick C. Beil, 1983. A wonderful (though expensive) compendium of typefaces,
  indexed in such a way that you can look a typeface up by its features, thus allowing you to
  home in quickly on an unknown specimen instead of spending hours leafing through
  catalogues. Lovers of my "horizontal and vertical problems" will delight in this book.
Phillips, Tom. A Humument: A Treated Victorian Novel. New York: Thames \& Hudson,
  1980. Like a child who has covered a wall with colorful crayon drawings, Tom Phillips has
  completely obliterated the pages of an old novel (A Human Document by W. H. Mallock)
  with his colorful scribblings, and only here and there do traces of the original show through.
  A droll stunt!
Poincare, Henri. "On Mathematical Creation". In The World of Mathematics, vol. 4, edited by
  James R. Newman, 2041-50. New York: Simon \& Schuster, 1956..A lecture presented to
  the Psychological Society in Paris in the early part of this century. Anticipating
  developments in cognitive science some eighty years later, Poincare speculates on the
  nature of the events taking place inside his skull as he makes mathematical discoveries.


Po1ya, George. How to Solve It. New York: Doubleday, 1957. P61ya, like Poincart- a
 mathematician fascinated by thought processes, attempts here to give recipes for how to
 attack mathematical problems. The problem with this is that there is-and can be-no failsafe
 recipe. Even trying to give guidelines is probably futile. The nose that smells the right route
 is simply rare, and there are no two ways about it.
Post, Emil. "Absolutely Unsolvable Problems and Relatively Undecidable Propositions:
 Account of an Anticipation". In The Undecidable, edited by Martin Davis, 338-443. New
 York: Raven, 1965. In this paper, Post concludes that mathematicians' thought processes are
 essentially creative and non-mechanizable. Related to the article by Myhill (see above).
Poundstone, William. The Recursive Universe: Cosmic Complexity and the Limits of
 Scientific Knowledge. New York: William Morrow, 1985. A superlative account of the
 reductionist miracle: fantastically complex entities-living, self-reproducing organisms-turn
 out to be vast arrays of very simply interacting parts. Von Neumann's trick of self-
 reproduction without infinite regress (adapted from Gödel) is one of the main topics
 explored here, and with the help of Conway's absorbing game of Life, which plays the
 starring role in his book, Poundstone does a masterful job of explaining how it comes about.
Racter. The Policeman's Beard Is Half Constructed. New York: Warner Books, 1984.
 "Racter" is a program written by Bill Chamberlain and Thomas Etter. The Policeman's
 Beard is a book written by Racter. It is all somewhat tongue-in-cheek, because Racter does
 not really know much about half-constructed beards, but what is lovely about Racier's prose
 is the way it skirts the fringes of meaning, weaving drunkenly across the boundary between
 sense and senselessness.
Rapoport, Anatol. Two-Person Game Theory. Ann Arbor: University of Michigan Press, Ann
 Arbor Science Library, 1966. A sound treatment of game theory, featuring a most
 .interesting personal discussion on opinions about the meaning of "rationality" in Prisoner's-
 Dilemma-like situations.
Reps, Paul. Zen Flesh, Zen Bones. New York: Doubleday, Anchor Press, n.d. An easily
 available collection of Zen koans, highly amusing and, perhaps, even enlighteningproviding
 you take it all with sufficiently many grains of salt.
Rogers, Hartley. Theory of Recursive Functions and Effective Computability. New York:
 McGraw-Hill, 1967. A standard reference work on many advanced concepts in
 metamathematics, including such concepts as "productive" and "creative" sets, referred to in
 my Chapter 13.
Ross, Alf. "On Self-Reference and a Puzzle in Constitutional Law". Mind 78, no. 309
 (January 1969): 1-24. A condundrum in the philosophy of law: Can laws modify
 themselves, or is that paradoxical? Ross' view is that logical inconsistency is unacceptable
 in law, and therefore that self-amendment is impossible. For a response, see the paper by
 Hart (above).
Rucker, Rudy. Infinity and the Mind. Boston: Birkhauser, 1982. A book that had to be
 written. Presenting the most abstruse concoctions of the mind in language that is not
 abstruse, and connecting it with thoughts about consciousness and the mystery of existence
 -this is what Rucker excels in.
Ruelle, David. "Les Attracteurs Etranges". La Recherche 11, no. 108 (February 1980): 131-
 44. A good article relating these wispy mathematical clouds to the physics they are
 supposed to explain, featuring a number of excellent illustrations.
Rumelhart, David E. and Donald A. Norman. "Simulating a Skilled Typist: A Study of
 Skilled Cognitive-Motor Performance". Cognitive Science 6, no. 3 (JulySeptember 1982):
 1-36. Anticipating the current "parallel distributed processing" project at the University of
 California at San Diego, the research described in this article is among the most interesting
 work on modeling human performance that I have encountered.
Russett, Bruce. The Prisoners of Insecurity. New York: W.H. Freeman, 1983. The arms race
 as an iterated Prisoner's Dilemma, and how we might be able to break out of the deadlock.
 The last section-"Responsibility"--concludes with this admonition (and how I wish people
 would take it to heart!): "In a democracy, silence about nuclear issues carries an implication
 not just of indifference but of acceptance. If we stand silent in the face of an arms race-and


  the war to which it may lead us-we must share responsibility for the outcome. `Silence
  gives consent.' "
Ryder, Frederick, and Company. Ryder Types, 2 vols. (with periodic supplements). Chicago:
  Frederick Ryder and Company. The best catalogue of typefaces I have run across-but it is
  expensive. Some of the oddest faces I have ever seen are found in the four supplements I
  own.
Sagan, Carl, ed. Communication with Extraterrestrial intelligence. Cambridge, Mass.: MIT
  Press, 1973.. An entertaining transcript of an international meeting in the days when
  Russians and Americans spoke to each other. Sagan puts forth his notion of various types of
  earth-life-based "chauvinisms". A hopeful and intellectually refreshing book-the kind one
  wishes there were hundreds of.
Sampson, Geoffrey. "Is Roman Type an Open-Ended System? A Response to ,Douglas
  Hofstadter". Visible Language 17, no. 4 (Autumn 1983): 410-12. Sampson tlisputes my
  claim that it is impossible to capture the fluid spirit of letters of the alphabet in parametrized
  computer subroutines; he suggests that it is possible, if you limit your goals to capturing the
  spirit of letters suitable for printing serious books in.
Schank, Roger. Dynamic Memory: A Theory of Reminding and Learning in Computers and
  People. New York: Cambridge University Press, 1982. Although I disagree with some of
  his theorizing, I agree fully with Schank's focus on the types of problems that cognitive
  science ought to be most concerned with, and I like the examples he uses.
Scherlis, William L. and Pierre L. Wolper. "Self-Referenced Referenced, and Self-
  Referenced". Communications of the ACM 23, no. 12 (December 1980): 736. A humorous
  short note on papers that cite themselves.
Schrodinger, Ernst. What Is Life? and Mind and Matter. New York: Cambridge University
  Press, 1967 (reprint of 1944 edition). This philosophically-minded physicist prophetically
  speculates on the nature of the hereditary message, before the days when DNA's structure or
  function were known. Also venturing into the deep waters of consciousness, he comes up
  with this mystical conclusion: "The over-all number of minds is just one."
Schwenk, Theodor. Sensitive Chaos. New York: Schocken, 1976. A book about fluids in the
  wild and in the laboratory, filled with striking patterns and fantastic photographs strongly
  suggesting this mystical conclusion: "The over-all number of fluids is just one."
Searle, John. "Minds, Brains, and Programs". The Behavioral and Brain Sciences 3
  (September 1980): 417-57. Like the Lucas article (see above), this one launched a thousand
  rebuttals (including two by me). In this, its original setting, it was followed by nearly thirty
  rebuttals and counter-rebuttals. It is amusing and educational to read them all. I view this
  article as a litmus test, in the sense that someone convinced by Searle's imagery is almost
  sure to have a very negative opinion of AI.
-----The Myth of the Computer: An Exchange". New York Review of Books (June 24, 1982):
  56-57. Searle portrays Dennett and me as blind advocates of "strong AI"-the notion that "the
  appropriately programmed computer literally has a mind". He resents the fact that such
  views are "well financed and backed by prestigious teams of research workers", and tells
  how he is constantly working toward "the relentless exposure of its preposterousness".
-----The Myth of the Computer". New York Review of Books (April 29, 1982): 3-6. A rather
  negative review of The Mind's I by someone roundly criticized in the book. The one good
  thing Searle does in this review is to stress the central epistemological problem facing AI: to
  explain the nature of the fluid reference, or semanticity, that mental activity exhibits.
Serafini, Luigi. Codex Seraphinianus. Milan: Franco Maria Ricci, 1981. A complete pseudo-
  encyclopedia, in two volumes. The writing system, the page numbering, the weird
  diagrams, and especially the wonderful color illustrations are all products of the cryptic
  mind of Serafini, an Italian architect. Also available in one volume, and more
  inexpensively, from Abbeville Press in New York, N.Y.
Seuss, Dr. On Beyond Zebra. New York: Random House, 1955. Humorous ways of extending
  the alphabet beyond 'z': a metaphor for jumping out of the system (' jootsing").
Simon, Herbert A. "Cognitive Science: The Newest Science of the Artificial". Cognitive
  Science 4, no. 2 (April June 1980): 33-46. Simon's confident assertion that computers


  already have what it takes to possess full intelligence: symbol manipulation. In his
  conclusion, he claims: "Wherever the boundary is drawn, there exists today a science of
  intelligent systems that extends beyond the limits of any single species."
-----Studying Human Intelligence by Creating Artificial Intelligence". American Scientist 69,
  no. 3 (May June 1981): 300-309. This is the lecture in which Simon spoke of the all-
  important 100-millisecond barrier, below which it is of no interest to cognitive scientists to
  know what happens (see my Chapter 26). In print he changed "100 millisecondsto "ten
  milliseconds", although the point remains the same either way.
Singmaster, David. Notes on Rubik's "Magic Cube". Hillside, NJ.: Enslow, 1981. On its
  cover, this book proudly boasts: "'The definitive treatise.'-Scientific American. " Well, if
  Scientific American says so, who am I to dispute it?
Sloman, Aaron. The Computer Revolution in Philosophy: Philosophy, Science, and Models of
  Mind. Atlantic Highlands, NJ.: Humanities Press, 1978. A book that warns philosophers
  that they will miss the boat if they don't jump on the Al bandwagon. This a good though
  somewhat tendentious discussion of the philosophical import of Al.
Smith, Brian C. "Reflection and Semantics in Lisp". Palo Alto, Calif.: Xerox Palo Alto
  Research Center Report, 1983. A boiling-down of a 500-page Ph.D. thesis into a dozen
  pages or so, concerning a system capable of reasoning about itself (and reasoning about
  such reasoning, etc.). This work exemplifies the "meta-meta" style of Al research (see my
  Chapter 23-especially its P.S. ).
Smith, Stephen B. The Great Mental Calculators: The Psychology, Methods, and Lives of
  Calculating Prodigies. New York: Columbia University Press, 1983. A set of portraits of
  very strange yet very human beings whose aberrant minds let them do calculating tasks that
  we ordinary people-no matter how number-loving-could not conceivably do.
Smolensky, Paul. "Harmony Theory: A Mathematical Framework for Stochastic Parallel
  Processing". University of California at San Diego Institute for Cognitive Science Technical
  Report ICS No. 8306. Based on the ideas of statistical mechanics, this project utilizes
  stochastic parallelism, regulated by a "temperature" that gradually drops to zero, to search
  most efficiently for the optimal global state of a system.
Smullyan, Raymond. This Book Needs No Title: A Budget of Living Paradoxes. Englewood
  Cliffs, NJ.: Prentice-Hall, 1980. A very humorous collection of observations about the
  constant intermingling of life and paradox, by a logician whose awareness of paradox is
  especially keen.
Solo, Dan X. Sans Serif Display Alphabets. New York: Dover, 1979. A collection of elegant
  letterforms whose subdued flair resides exclusively in their gentle curves, stroke taperings,
  line endings, and the interplay between positive and negative space.
-----Special-Effects and Topical Alphabets. New York: Dover, 1978. After you've looked at a
  book like this, you know why the problem of letterforms is synonymous with the problem
  of full human intelligence.
Sonneborn, Tracy M. "Degeneracy of the Genetic Code: Extent, Nature, and Genetic
  Implications". In Evolving Genes and Proteins, edited by Vernon Bryson and
Henry J. Vogel. New York: Academic Press, 1965. Perhaps the first paper to suggest that
  there is evolutionary rhyme and reason to the particular match-up between codons and
  amino acids that exists in the arbitrary-seeming genetic code.
Soppeland, Mark. Words. Los Altos, Calif.: William Kaufmann, 1980. A witty collection of
  words drawn as self-referential pictures: stacks of books whose shapes spell out "books",
  and so on.
Sorrels, Bobbye. The Nonsexist Communicator. Englewood Cliffs, NJ.: Prentice-Hall, 1983.
  A sizable collection of sexist usages and nonsexist remedies; some of the remedies,
  however, are needlessly awkward.
Sperry, Roger. "Mind, Brain, and Humanist Values". In New Views on the Nature of Man,
  edited by John R. Platt. Chicago: University of Chicago Press, 1965. The article that asks,
  and tries to answer, the question "Who pushes whom around in the population of causal
  forces that occupy the cranium?" My answer is in Chapter 25.



Spinelli, Aldo. Loopings. Amsterdam: Multi-Art Points Edition, 1976. A somewhat clumsy
  realization of a good idea: a book whose pages are all self-inventorying sentences, or close
  to that. At the end, Spinelli discusses the kinds of locked-in loops and entryways into
  looping that can arise.
Stanley, H. Eugene et al. "Interpretation of the Unusual Behavior of H2O and D20 at Low
  Temperature: Are Concepts of Percolation Relevant to the 'Puzzle of Liquid Water'?"
  Physica. 106a (1981): 260-77. An article that gives a fairly recent picture of the "flickering-
  cluster" view of water.
Stein, Gertrude. How to I Write. Craftsbury Common, Vt.: Sherry Urie, 1977. (Originally
  published in 1931.) Classic volume of absurdities. It would be a nightmare to try to read the
  whole thing; each page is like a unique grain of sand, but all the pages taken together add up
  to a rather bland beach. Better to savor it in small bits.
Steiner, George. After Babel: Aspects of Language and Translation. New York: Oxford
  University Press, 1975. This fascinating exploration of the depths of personal meaning is
  full of obscurities, but it is still one of the best statements on the subtlety of language that 'I
  know.
Strich, Christian. Fellini's Faces: 418 Photographs from the Archives of Federico Fellini. New
  York: Holt, Rinehart, and Winston, 1981. A celebration of the remarkable diversity of
  human physiognomies, as well as their amazing expressivity.
Stryer, Lubert. Biochemistry. New York: W. H. Freeman, 1975. An elegantly produced text
  of biochemistry, on a slightly lower level than Lehninger. The figures-many in color -are
  especially appealing.
Suber, Peter. "A Bibliography of Works on Reflexivity". Unpublished manuscript, 1984. An
  extensive collection of pointers to the vast universe of writings on self-modifying laws, self-
  referring literature, self-fulfilling prophecies, self-replicating machines, self-monitoring
  computer programs, and on and on.
-----The Paradox of Self-Amendment: A Study of Logic, Law, Omnipotence, and Change.
  Forthcoming. This is the first (and only) book-length study of logical paradoxes in law. It
  focuses on one family of paradoxes (rules that authorize change being applied to
  themselves), but eventually tries to address the general problem of how legal reasoning
  copes with logical paradox. The game of Nomic is featured in an appendix.
Thomas, Dylan. Collected Poems. New York: New Directions, 1953.1 am the first to admit I
  don't understand poetry well, but this beguiling collection, with its combination of beautiful
  language and bewildering opacity, leaves me especially disturbed.
Thomas, Lewis. The Medusa and the Snail. New York: Viking, 1979. Short essays on science
  and life, some of which are very insightful, others of which I find puzzling or just plain
  silly.
Turing, Alan. "Computing Machinery and Intelligence". Reprinted in Minds and Machines,
  edited by Alan Ross Anderson. Englewood Cliffs, NJ.: Prentice-Hall,
1964. In which the now-infamous "Turing Test" for machine thought is proposed. Lucid and
  straightforward, this article contains plenty of provocative ideas, even today.
Turner-Smith, Ronald. The Amazing Pyraminx. Hong Kong: Meffert Novelties, 1981. A
  simple guide to the notation, group theory, and solution of the Pyraminx puzzle.
Ulam, Stanislaw. Adventures of a Mathematician. New York: Scribners, 1976. A fascinating
  account of the intellectual life of a highly innovative and fun-loving mathematician,
  including some speculations on mind and consciousness.
Vetterling-Braggin, Mary, ed. Sexist Language: A Modern Philosophical Analysis. Totowa,
  NJ.: Littlefield, Adams \& Co., 1981. A valuable collection of articles from many
  perspectives on sexism. Includes one section on the comparison between sexism and
  racism.
von Neumann, John. Theory of Self-Reproducing Automata, edited and completed by
  Arthur W. Burks. Urbana, Illinois: University of Illinois Press, 1966. Here von Neumann
  describes how a machine could construct replicas of itself, or even machines more complex
  than itself, out of raw materials. At the core of this work by von Neumann, however, is



  Gödel’s method of achieving mathematical self-reference-a fact that many people seem to
  overlook or downplay.
Walker, Alan, ed. The Chopin Companion. New York: W. W. Norton, 1966. An excellent
  collection of articles by composers, performers, and musicologists on Chopin, his music,
  and its influence.
Watzlawick, Paul. Change. New York: W. W. Norton, 1974. A theory of psychological
  disorder as caused by encounters with paradox in daily living. The suggested therapy is to
  fight paradox with paradox (if that isn't paradoxical!).
Webb, Judson. Mechanism, Mentalism, and Metamathematics. Hingham, Mass.: D. Reidel,
  1980. A penetrating scholarly analysis of the import of the most important
  metamathematical work by Gddel, Church, Turing, Kleene, Tarski, Post, and others. In
  particular, it seeks to elucidate the relationship of the famous "limitative theorems"
  applicable to formal systems with attempts to mechanize human mental activity. In essence,
  its conclusion is that those results strengthen, rather than diminish, the notion that mental
  activity can in principle be mechanized.
Wells, Carolyn. A Nonsense Anthology. New York: Dover Press, 1958. A wonderful
  collection of strange poems and silly snippets of pseudo-sense.
Wheelis, Allen. The Scheme of Things. New York: Harcourt Brace Jovanovich, 1980. A
  moving story of human and animal attachments, and the emotional resonances of places and
  scenes. The story hints at many self-references, most notably to another novel called The
  Way Things Are, written by Wheelis' protagonist Oliver Thompson who, like Wheelis
  himself, is an unorthodox psychiatrist living in San Francisco. Wheelis within Wheelis!
Williams, J. D. The Compleat Strategyst. New York: McGraw-Hill, 1954. One of the earliest
  and clearest books ever written on game theory. It is a very elementary book, with
  humorous drawings and droll stories to illustrate its points.
Wilson, E. O. The Insect Societies. Cambridge, Mass.: Harvard University, Belknap Press,
  1971. It is heavy going to read this book, but parts of it are engrossing, as they tell how
  order at a high level comes from the independent activity of low-level organisms acting
  without knowledge of one another.
Winston, Patrick Henry. The Psychology of Computer Vision. New York: McGraw-Hill,
  1975. Although fairly old by AI standards, this collection of six articles is still of
  considerable interest. One is Minsky's article on frames; another is an article by Winston on
  learning and recognition; and another is an article by Waltz on economical strategies for
  vision in the "blocks world".
Winston, Patrick Henry and Berthold Horn. Lisp. Reading, Mass.: Addison-Wesley, . 1981. A
  solid textbook on Lisp, starting right at' the beginning and going up to AI programming
  techniques.
Yaguello, Marina. Alice au pays du langage. Paris: Editions du Seuil, 1981. A French linguist
  examines the phonetics, syntax, semantics, and gradual shifts of her own language to
  illuminate the strangeness of human thought. Many of the examples in the book are taken
  from popular language, slang, or humor, which makes it particularly unstuffy.
-----Les mots et les femmes. Paris: Petite Bibliotheque Payot, 1978. On fighting the battle
  against sexist language in French-but unfortunately, because of the two-gender system of
  romance languages, eradication of sexism will prove far more difficult in those languages
  than in English (where it is already ferociously hard).
Zapf, Hermann. About Alphabets. Cambridge, Mass.: MIT Press, 1960. A famous
  contemporary type designer describes his adventures among the alphabets of the world.



